\section{Concorrência Estruturada}
\label{sec:structured}

\subsection{Implementação em Java}

\subsubsection{Descrição}

A v4 parte da v3 --- a variante com melhor desempenho até então --- e não
altera o algoritmo nem o padrão de acumulação local por tarefa. A mudança
é inteiramente na camada de \emph{coordenação}: o par
\texttt{ExecutorService} + \texttt{CountDownLatch} da v3 é substituído por
\texttt{StructuredTaskScope} (JEP 480) e o snapshot de centróides, antes
passado como parâmetro através de \texttt{processChunk}, passa a ser
propagado via \texttt{ScopedValue} (JEP 446). Ambas são APIs de
\textit{preview} no JDK 25, exigindo a flag \texttt{--enable-preview} em
toda invocação da JVM que carregue \texttt{KMeans.class} --- compilação,
JMH, JCStress, JMeter e execução direta.

Essa troca não é uma otimização de desempenho: é uma correção de
robustez. Na v3, uma exceção lançada dentro de uma tarefa de worker não
tinha canal de propagação de volta à thread coordenadora --- o
\texttt{CountDownLatch} apenas conta conclusões, sem carregar erros. Na
prática, isso significa que uma falha real ficaria silenciosamente
mascarada, manifestando-se apenas de forma indireta mais adiante (por
exemplo, como um resultado ausente na redução). O
\texttt{StructuredTaskScope} resolve isso estruturalmente:
\texttt{scope.join()} propaga a exceção original de qualquer subtarefa
para a thread coordenadora e cancela automaticamente as subtarefas
irmãs ainda em execução --- sem necessidade de código adicional de
tratamento de erro.

\begin{table}[H]
\centering
\begin{tabular}{lp{11cm}}
\hline
\textbf{Componente} & \textbf{Alteração} \\
\hline
\texttt{pom.xml}    & flag \texttt{--enable-preview} adicionada ao \texttt{maven-compiler-plugin} \\
\texttt{KMeans.java} & \texttt{ExecutorService} + \texttt{CountDownLatch} substituídos por
                       \texttt{StructuredTaskScope}; centróides propagados por
                       \texttt{ScopedValue<double[][]> CENTROIDS} em vez de parâmetro de método \\
\texttt{justfile}   & flag \texttt{--enable-preview} adicionada a todas as invocações de
                       \texttt{java}/JMH/JCStress/JMeter \\
\hline
\end{tabular}
\caption{Alterações da v4 em relação à v3.}
\label{tab:v4-changes}
\end{table}

O confinamento local por tarefa (\texttt{localSums}/\texttt{localCounts}
dentro de \texttt{processChunk}) é mantido exatamente como na v2/v3: a
v4 contribui apenas na camada de coordenação, não na camada de
confinamento, que já havia sido validada por JCStress nas versões
anteriores.

\subsubsection{Resultados}

\paragraph{Avaliação de Condição de Corrida}

Os mesmos quatro testes JCStress (JCS1--JCS4) das seções anteriores
foram executados sobre a v4, com a mesma configuração (preset
\texttt{default}, 6 forks, 5 iterações de 1.000 ms por fork).

\begin{table}[H]
\centering
\begin{tabular}{llll}
\hline
\textbf{ID} & \textbf{Resultado} & \textbf{Outcome FORBIDDEN} & \textbf{Freq.\ máx.} \\
\hline
JCS1 & \textbf{PASSOU}            & nenhum                   & ---      \\
JCS2 & \textbf{PASSOU}            & nenhum                   & ---      \\
JCS3 & FALHOU (esperado) & \texttt{"1"} lost update em \texttt{counts[0]++}  & 30,00\% \\
JCS4 & FALHOU (esperado) & \texttt{"1.0"} lost update em \texttt{sums[0]+=1.0} & 29,61\% \\
\hline
\end{tabular}
\caption{Resultados JCStress --- regiões críticas da v4.}
\label{tab:jcstress-v4}
\end{table}

JCS1 e JCS2 continuam passando: a leitura concorrente do centróide via
\texttt{ScopedValue} é segura, assim como a leitura direta por parâmetro
nas versões anteriores. JCS3 e JCS4 falham como esperado --- o teste
exercita deliberadamente o mesmo padrão de array compartilhado sem
proteção, independente da camada de coordenação usada. A frequência
máxima observada (30,00\% e 29,61\%) está na mesma ordem de grandeza da
v2 (28,31\%/27,10\%); a frequência exata de um outcome de \textit{lost
update} é inerentemente dependente do escalonamento de threads no
momento da execução, não um valor que deveria se repetir entre execuções
da mesma condição de corrida.

\paragraph{Avaliação com Microbenchmark}

Os microbenchmarks foram executados com JMH 1.37 sobre JDK 25.0.1, com
\texttt{--enable-preview}, 1 fork, 1 iteração de aquecimento de 10 s e 3
iterações de medição de 10 s cada.

\begin{table}[H]
\centering
\begin{tabular}{lp{4.2cm}rrrr}
\hline
\textbf{ID} & \textbf{Benchmark} & \textbf{v3} & \textbf{v4} & \textbf{Erro v4 ($\pm$)} & \textbf{Un.} \\
\hline
B1 & Execução completa (\textit{k}=10, iter=20) & 4{,}55 & 7{,}08  & 1{,}26 & s/op \\
B2 & Por iteração (\textit{k}=10, iter=5)        & 1{,}60 & 1{,}98  & 0{,}07 & s/op \\
B3 & \texttt{closestCentroid} (38 dims)           & 96{,}82 & 131{,}18 & 97{,}71 & ns/op \\
B4 & \texttt{squaredDistance} (38 dims)           &  9{,}31 &  12{,}89 &  6{,}79 & ns/op \\
\hline
\end{tabular}
\caption{Resultados JMH --- v3 vs.\ v4.}
\label{tab:jmh-v4}
\end{table}

Todos os quatro benchmarks pioram levemente em relação à v3: B2 sobe de
1,60 s para 1,98 s (\textasciitilde24\% mais lento), e B3/B4 --- que
medem exatamente o mesmo código aritmético de \texttt{closestCentroid} e
\texttt{squaredDistance}, inalterado desde a v1 --- também aumentam,
embora com erro elevado. Como o algoritmo em si não mudou, esse custo
adicional é atribuído à sobrecarga das APIs de \textit{preview}: tanto
\texttt{StructuredTaskScope} quanto \texttt{ScopedValue} ainda não têm a
maturidade de compilação JIT do caminho
\texttt{ExecutorService}+\texttt{CountDownLatch}, que é parte estável da
plataforma há muito mais tempo. Este é um resultado honesto a relatar,
não um problema a esconder: a v4 troca uma pequena perda de desempenho
por uma coordenação estruturalmente mais robusta.

\paragraph{Avaliação com Macrobenchmark}

Os macrobenchmarks foram executados com JMeter 5.6.3 ($k=10$,
\texttt{maxIter}=3, \texttt{tol}=0,0), 120 s por configuração e ramp-up
de 30 s. Apenas o coletor \textbf{ParallelGC} foi testado --- decisão de
escopo para a v4 (e v5), sem repetir a varredura de três coletores feita
em v1--v3.

\begin{table}[H]
\centering
\begin{tabular}{rrrr}
\hline
\textbf{Threads} & \textbf{Execuções} & \textbf{Tempo (s)} & \textbf{Throughput (exec/min)} \\
\hline
1  &  92 & 120{,}4 & 45{,}8 \\
2  &  97 & 121{,}3 & 48{,}0 \\
4  &  99 & 123{,}0 & 48{,}3 \\
8  & 101 & 124{,}4 & 48{,}7 \\
16 & 115 & 132{,}7 & 52{,}0 \\
\hline
\end{tabular}
\caption{Throughput (execuções/min, ParallelGC) por número de threads JMeter --- v4.}
\label{tab:jmeter-v4}
\end{table}

Comparado à v3 com o mesmo coletor (52,5 a 62,3 exec/min de 1 a 16
threads), a v4 fica abaixo em todos os níveis de carga (45,8 a 52,0
exec/min) --- consistente com o aumento já observado em B2: o custo
adicional por iteração se propaga para o throughput agregado.

\paragraph{Avaliação de Profile}

O profile foi capturado com JFR (\texttt{settings=profile}) sobre uma
execução completa do algoritmo. A gravação totalizou 1.806 amostras de
\texttt{jdk.ExecutionSample}.

\begin{table}[H]
\centering
\small
\begin{tabular}{p{7.4cm}rrp{3cm}}
\hline
\textbf{Método} & \textbf{Amostras} & \textbf{CPU (\%)} & \textbf{Categoria} \\
\hline
\texttt{KMeans.squaredDistance}       & 588 & 32{,}6 & Algoritmo \\
\texttt{String.split}                 & 181 & 10{,}0 & Parsing CSV (inicialização) \\
\texttt{FloatingDecimal.copyDigits}   & 162 &  9{,}0 & Parsing CSV (inicialização) \\
\texttt{Arrays.copyOf}                & 159 &  8{,}8 & Coordenação (subtarefas) \\
\texttt{KMeans.closestCentroid}       & 157 &  8{,}7 & Algoritmo \\
\texttt{FloatingDecimal.readJavaFormatString} & 125 &  6{,}9 & Parsing CSV (inicialização) \\
\texttt{KMeans.addPoint}              &  92 &  5{,}1 & Algoritmo \\
\texttt{FloatingDecimal.doubleValue}  &  75 &  4{,}2 & Parsing CSV (inicialização) \\
\texttt{BufferedReader.readLine}      &  26 &  1{,}4 & I/O \\
\hline
\end{tabular}
\caption{Hot methods na v4 (JFR, 1.806 amostras).}
\label{tab:jfr-hot-v4}
\end{table}

O algoritmo (\texttt{squaredDistance} + \texttt{closestCentroid} +
\texttt{addPoint}) soma \textasciitilde46\% das amostras, o maior bloco
individual. Chama atenção a persistência de frames de \textit{parsing}
de CSV (\texttt{String.split}, \texttt{FloatingDecimal.*}) somando
\textasciitilde39\%: isso não é um retorno do gargalo de I/O por
iteração eliminado na v3 --- é a leitura única do dataset
(\texttt{loadAllPoints}, executada uma vez antes do laço de iterações)
caindo dentro da janela de gravação do JFR, mesmo estando excluída da
medição de tempo do JMH/JMeter. O bloco \texttt{Arrays.copyOf}
(\textasciitilde9\%) é overhead de coordenação --- a lista de
\texttt{StructuredTaskScope.Subtask} sendo montada e convertida a cada
iteração.

Foram registradas 61 coletas de GC, com pausa total de 1.013,97 ms e
pausa média de 16,62 ms por evento. O heap atingiu pico de 2.764,8 MB
(\textasciitilde2,76 GB), em linha com o \textasciitilde2,7 GB da v3 ---
esperado, já que ambas mantêm o mesmo dataset em memória como
\texttt{double[][]}. GC não é gargalo em nenhuma das duas versões.

\subsection{Discussão}

A Tabela~\ref{tab:jmh-comparativo-v4} resume a evolução do custo por
iteração desde a versão serial.

\begin{table}[H]
\centering
\begin{tabular}{lrrrr}
\hline
\textbf{Abordagem} & \textbf{B1 (s/op)} & \textbf{B2 (s/op)} & \textbf{Erro B1} & \textbf{Erro B2} \\
\hline
Serial & 212{,}57 & 60{,}05 & $\pm$34{,}48 & $\pm$2{,}46 \\
v3     &   4{,}55 &  1{,}60 & $\pm$17{,}98 & $\pm$0{,}17 \\
v4     &   7{,}08 &  1{,}98 &  $\pm$1{,}26 & $\pm$0{,}07 \\
\hline
\end{tabular}
\caption{Comparativo JMH: Serial, v3 e v4 (s/op).}
\label{tab:jmh-comparativo-v4}
\end{table}

A v4 não é uma otimização de desempenho --- é uma atualização de idioma
de concorrência. A v3 tinha uma lacuna estrutural real: por usar
\texttt{ExecutorService}+\texttt{CountDownLatch}, uma exceção lançada
dentro de uma tarefa de worker não possuía canal de propagação de volta
à thread coordenadora, podendo ficar silenciosamente mascarada. O
\texttt{StructuredTaskScope} fecha essa lacuna por construção:
\texttt{scope.join()} propaga a exceção original e cancela
automaticamente as subtarefas irmãs. O \texttt{ScopedValue} substitui a
passagem do snapshot de centróides por parâmetro de método, sendo a
alternativa moderna ao \texttt{ThreadLocal} --- que não se propaga
corretamente para virtual threads filhas criadas dentro do escopo.

O preço dessa robustez é um aumento modesto no tempo por iteração (B2:
1,60 s $\to$ 1,98 s) e no throughput agregado (52,5--62,3 $\to$
45,8--52,0 exec/min), atribuído à imaturidade de compilação das duas
APIs de \textit{preview}. O JCStress confirma que o confinamento local
por tarefa --- a garantia que realmente importa para corretude ---
permanece inalterado: JCS1/JCS2 continuam passando, e JCS3/JCS4
continuam falhando exatamente como esperado quando a proteção é
removida deliberadamente.
